% Revised 30 August 2026
\documentclass[10pt]{article}
\usepackage[margin=0.75in]{geometry}
\usepackage[T1]{fontenc}
\usepackage{lmodern,microtype}
\usepackage{amsmath,amssymb,amsthm,mathtools}
\usepackage{enumitem,booktabs,tabularx,xcolor,hyperref}
\hypersetup{colorlinks=true,linkcolor=blue!45!black,urlcolor=blue!45!black}
\setlength{\parindent}{0pt}
\setlength{\parskip}{4pt}
\setlist{leftmargin=*,itemsep=1pt,topsep=2pt}
\newcommand{\Z}{\mathbb Z}
\newcommand{\Gspace}{\mathfrak G}
\newtheorem{theorem}{Theorem}

\begin{document}
\begin{center}
{\LARGE\bfseries Research Statement\par}
\vspace{1mm}
{\large Local information and global structure in countable graphs\par}
\vspace{2mm}
{Hadi Zamani --- Computer Science Institute, Charles University\par}
{Supervisor: Robert \v{S}\'amal\par}
{\small 30 August 2026}
\end{center}

\begin{abstract}
My research concerns local-to-global questions in sparse and countable graphs.  My
current joint work with Robert \v{S}\'amal has three concrete parts: extending
nowhere-zero flow theory to countably infinite locally finite graphs, comparing
Baire-category notions of typicality in rooted and unrooted bounded-degree graph
spaces, and proving an undecidability theorem for reconstruction from exact finite
local data.  A further direction is to understand the recent explicit construction of
a nonsofic group and to investigate whether its non-approximability argument has a
direct combinatorial formulation.
\end{abstract}

\section{Guiding question}

A bounded-degree graph is observed locally through its finite rooted balls, but global
features such as ends, bridges, flows, and approximation by finite graphs need not be
determined by any one ball.  My guiding question is:
\begin{quote}
\emph{Which global properties of a countable locally finite graph are forced by its
finite neighbourhoods, and how does the answer depend on what local information is
retained?}
\end{quote}
The last clause is essential.  A rooted topology records the neighbourhood of one
distinguished vertex; an unrooted local-spectrum topology records which rooted
neighbourhood types occur anywhere; Benjamini--Schramm convergence records their
frequencies.  These observation models answer different questions and should not be
identified.

\section{Established work: flows and bridges}

My current joint manuscript with Robert \v{S}\'amal is entitled
\emph{Typical properties of countable graphs: flows and bridges}.  Its first part
studies nowhere-zero flows in countably infinite locally finite graphs.

Fix an orientation of a locally finite graph $G$.  For an abelian group $A$, a
\emph{weak $A$-flow} assigns an element of $A$ to every directed edge and satisfies
Kirchhoff's law at every vertex.  A \emph{strong $A$-flow} requires the corresponding
identity across every finite edge cut.  These conditions coincide for finite graphs,
but not for infinite graphs.

If $e$ is a bridge, then $G-e$ has two components.  We call $e$ \emph{one-way
infinite} if one component is finite, and \emph{two-way infinite} if both components
are infinite.  Summing the vertex equations over the finite component shows that a
weak flow must vanish on every one-way infinite bridge.  Our theorem proves that this
necessary obstruction is also sufficient for a nowhere-zero $\Z_6$-flow.

\begin{theorem}[Weak infinite $6$-flow theorem]
Let $G$ be connected, countably infinite, and locally finite.  Then $G$ admits a weak
nowhere-zero $\Z_6$-flow if and only if $G$ has no one-way infinite bridge.
\end{theorem}

For sufficiency, root $G$ and form finite multigraphs by contracting the complement
of successively larger rooted balls.  The absence of a one-way infinite bridge makes
the relevant finite contractions bridgeless.  Seymour's finite $6$-flow theorem gives
a nowhere-zero $\Z_6$-flow on each finite model, and compactness produces a single
edge assignment satisfying all vertex equations of $G$.  The proof therefore depends
on choosing finite models that preserve exactly the obstruction relevant to weak
conservation.

\section{Established work: Baire-category typicality}

There is no canonical uniform probability distribution on isomorphism classes of all
countably infinite bounded-degree graphs.  In the spaces considered in our manuscript,
we instead call a property \emph{typical} when its class is comeagre.  The definition is
always relative to a specified topology.

In the rooted space $\Gspace^{\bullet}_{\leq D}$, two graphs are close when large
balls around their roots are isomorphic.  For every $D\geq3$, the class of bridgeless
graphs is nowhere dense.  The analogous statement in the rooted $D$-regular space is
proved for odd $D$.  In the rooted spaces treated in the manuscript, graphs with a
two-way infinite bridge form a meagre class, while graphs with no cut vertex and
$D$-edge-colourable graphs form nowhere-dense classes for $D\geq3$.  The proofs extend
an arbitrary prescribed rooted ball far from the root so that every graph in a smaller
open neighbourhood contains the required obstruction.

The unrooted setting uses occurrence rather than one distinguished viewpoint.  Define
the radius-$r$ local spectrum by
\[
  \operatorname{Spec}_r(G)
  =\bigl\{[B_G(v,r),v]:v\in V(G)\bigr\},
\]
with multiplicities ignored.  Comparing these finite sets for increasing $r$ gives a
pseudometric: nonisomorphic graphs may have distance zero when they have exactly the
same finite rooted neighbourhood types.

We construct infinite \emph{strongly isolated} graphs, meaning that some finite radius
separates the graph from every nonisomorphic graph in the space.  These examples give
the following nonmeagreness statements in $\Gspace^{\infty}_{\leq D}$:
\begin{enumerate}
\item for every $D\geq3$, the graphs with a one-way infinite bridge and the graphs
      without one are both nonmeagre;
\item for every $D\geq3$, the graphs with a two-way but no one-way infinite bridge
      form a nonmeagre class;
\item for every $D\geq4$, the bridgeless graphs form a nonmeagre class.
\end{enumerate}
Thus the rooted and unrooted conclusions are genuinely different.  This contrast is
one of the main results of the manuscript; it is not a change of terminology, but a
change in the information encoded by the topology.

\section{Established work: undecidability from exact local data}

The manuscript also studies whether finite local information determines a unique
infinite graph.  Given a radius $r$ and a finite family $\mathcal F$ of rooted
radius-$r$ graphs, let
\[
 \Gspace_{r,\mathcal F}^{=,\infty}
 :=\bigl\{[G]:G\text{ is countably infinite, connected, locally finite, and }
              \operatorname{Spec}_r(G)=\mathcal F\bigr\}.
\]
The existence problem asks whether this class is nonempty.  Our new result concerns
uniqueness after existence has been promised.

\begin{theorem}[Promised uniqueness is undecidable]
Even under the promise that
$\Gspace_{r,\mathcal F}^{=,\infty}\neq\varnothing$, there is no algorithm that
always decides whether
\[
  \bigl|\Gspace_{r,\mathcal F}^{=,\infty}\bigr|=1.
\]
The statement remains true when every realization is required to be a simple
unlabelled graph of maximum degree at most $3$.
\end{theorem}

Equivalently, exact finite local data may determine one isomorphism type or may admit
at least two nonisomorphic infinite realizations, but no algorithm always distinguishes
these cases, even under the nonemptiness promise.  The reduction passes from finite
presentations to connected covers and then to exact-spectrum graphs; the subcubic
encoding shows that the obstruction is not caused by labels, multiple edges, or
unbounded degree.

\section{A direction for study: the explicit nonsofic construction}

The Aldous--Lyons conjecture asked whether every unimodular random rooted network is
a local weak limit of finite networks.  It was disproved by Bowen, Chapman, Lubotzky,
and Vidick, together with Bowen, Chapman, and Vidick~\cite{BCLV,BCV}.  That result
established the existence of non-approximable unimodular networks, but did not produce
a nonsofic group.

The 2026 OpenAI report \emph{Ten Advances in Mathematics and Theoretical Computer
Science} gives an explicit nonsofic group~\cite{OAI}.  More precisely, with
$R=L_{\mathbb F_2}(1,2)$, its proof constructs the finitely generated subgroup
\[
  G\cong \operatorname{EL}_9(R)\leq R^{\times}
\]
and proves that $G$ is not sofic.  For a finite generating tuple, the corresponding
directed generator-marked Cayley diagram is therefore a deterministic unimodular
network that is not a local weak limit of finite networks with the same markings.

The known argument proves non-approximability through soficity, group actions, and the
algebraic structure of the Leavitt group.  I would like first to understand this proof
and the role of each of these ingredients.  A natural question arising from that study
is:
\begin{quote}
\emph{Does there exist a direct combinatorial construction and proof of a
bounded-degree unimodular network that is not a Benjamini--Schramm limit of finite
graphs, without using group soficity or other group-theoretic machinery?}
\end{quote}
The immediate aim is not to claim a new construction.  It is to identify which steps
of the existing proof are genuinely group-theoretic and which can be expressed using
finite neighbourhoods, expansion, matchings, or counting.  This may clarify whether a
direct graph-theoretic proof is realistic and how the result fits the theory of local
convergence.

\section{Research plan}

My immediate programme is deliberately limited to questions for which the object and
the required conclusion are explicit:
\begin{enumerate}
\item \textbf{Complete the flows-and-typicality manuscript.}  Finalize the proofs and
      state every result with its degree bound, rooted or unrooted ambient space, and
      weak or strong conservation convention; include the exact-spectrum
      undecidability theorem with its nonemptiness promise and subcubic restriction.
\item \textbf{Study regular strong isolation.}  Determine whether infinite
      $D$-regular strongly isolated graphs exist and whether the degree threshold for
      the bridgeless unrooted construction can be reduced from four to three.
\item \textbf{Understand the explicit nonsofic construction.}  Study the OpenAI
      proof, identify its combinatorial core, and investigate whether any part of the
      non-approximability argument can be reformulated directly in the language of
      bounded-degree graphs.
\item \textbf{Compare approximation with category.}  Formulate the unlabelled
      bounded-degree problem in the space of unimodular laws and study the ergodic
      subspace separately from the full convex space.
\end{enumerate}

The flow, Baire-category, and promised-uniqueness theorems are results of the current
manuscript.  The regular-isolation, nonsofic-construction, and ergodic approximation
topics are directions for further study.

\textbf{Current research output.}
R.~\v{S}\'amal and H.~Zamani,
\emph{Typical properties of countable graphs: flows and bridges}, manuscript in
preparation (2026).

\renewcommand{\refname}{Selected references}
\begin{thebibliography}{9}
\footnotesize
\setlength{\itemsep}{0pt}
\setlength{\parskip}{0pt}
\bibitem{BCLV}
L.~Bowen, M.~Chapman, A.~Lubotzky, and T.~Vidick,
\emph{The Aldous--Lyons conjecture I: Subgroup tests},
arXiv:2408.00110, 2024.

\bibitem{BCV}
L.~Bowen, M.~Chapman, and T.~Vidick,
\emph{The Aldous--Lyons conjecture II: Undecidability},
arXiv:2501.00173, 2025.

\bibitem{OAI}
OpenAI, \emph{Nonsofic groups exist}, Chapter~3 in
\emph{Ten Advances in Mathematics and Theoretical Computer Science}, 2026,
\url{https://cdn.openai.com/pdf/ten-proofs-oai.pdf}.
\end{thebibliography}

\end{document}